\documentclass[11pt]{article}
\usepackage[margin=1in]{geometry}
\usepackage{amsmath, amssymb, amsthm}
\usepackage{algorithm}
\usepackage{algpseudocode}
\usepackage{graphicx}
\usepackage{hyperref}
\usepackage{booktabs}

\title{Guided Reverse Diffusion for Optimal Stimulus Generation}
\author{Internal Reference Document}
\date{March 2026}

\newcommand{\R}{\mathbb{R}}
\newcommand{\E}{\mathbb{E}}
\newcommand{\N}{\mathcal{N}}

\begin{document}
\maketitle

\begin{abstract}
We describe a method for generating natural-looking images that maximize the
information gained about a neuron's response properties, by combining a
denoising diffusion model with a Gaussian Process (GP) utility function.
The diffusion model acts as a learned prior over natural images, while the GP
utility gradient steers the generation toward informative stimuli.
This document is self-contained: it covers the diffusion model, the GP utility,
and how they are combined.
\end{abstract}

%% ============================================================
\section{Problem Setup}

We have a neuron whose firing rate $f(x)$ depends on a visual stimulus
$x \in \R^{d}$ (a $64 \times 64$ grayscale image, so $d = 4096$).
We model the log-firing rate via a Gaussian Process:
\[
  \lambda(x) \sim \mathcal{GP}(\mu(x), k(x, x')),
  \qquad
  f(x) = \exp\bigl(A \cdot \lambda(x) + \lambda_0\bigr),
\]
where $A$ and $\lambda_0$ are learned likelihood parameters.
Given a training set of stimulus--response pairs, we fit the GP and obtain a
posterior $q(\lambda \mid D)$.

\medskip
\noindent\textbf{Goal.}
Generate a new image $x^*$ that:
\begin{enumerate}
  \item Maximizes an information-theoretic utility $U(x^*)$ (how much we learn
        about the neuron by showing $x^*$), and
  \item Looks like a natural image (so it can be physically displayed on a
        projector and is ecologically meaningful).
\end{enumerate}

Requirement~1 alone can be solved by gradient ascent on $U$ in pixel space,
but the resulting images are unnatural (black-and-white noise patterns).
The diffusion model enforces Requirement~2 by constraining the generation
to the learned manifold of natural images.


%% ============================================================
\section{Background: Denoising Diffusion Models}

A Denoising Diffusion Probabilistic Model (DDPM) learns to reverse a gradual
noising process.  We summarize the key ideas; see Ho et al.\ (2020) for full
details.

\subsection{Forward Process}

Starting from a clean image $x_0$, the forward process adds Gaussian noise over
$T = 1000$ timesteps according to a variance schedule
$\beta_1, \ldots, \beta_T$:
\[
  q(x_t \mid x_0) = \N\!\bigl(x_t;\; \sqrt{\bar\alpha_t}\, x_0,\;
  (1 - \bar\alpha_t)\, I\bigr),
\]
where $\alpha_t = 1 - \beta_t$ and
$\bar\alpha_t = \prod_{s=1}^{t} \alpha_s$ is the cumulative signal retention.

\medskip
\noindent\textbf{Noise schedule.}
We use a linear schedule with $\beta_1 = 10^{-4}$ and $\beta_T = 0.02$.
At $t = 0$, $\bar\alpha_0 \approx 1$ (clean image);
at $t = T$, $\bar\alpha_T \approx 0$ (pure noise).

The forward process can be sampled in one step (``reparameterization trick''):
\begin{equation}\label{eq:forward}
  x_t = \sqrt{\bar\alpha_t}\, x_0 + \sqrt{1 - \bar\alpha_t}\, \varepsilon,
  \qquad \varepsilon \sim \N(0, I).
\end{equation}

\subsection{Reverse Process and the UNet}

A neural network $\varepsilon_\theta(x_t, t)$ --- a UNet with 99.5M parameters
--- is trained to predict the noise $\varepsilon$ added at step $t$.
Given a noisy image $x_t$ and the noise prediction $\hat\varepsilon = \varepsilon_\theta(x_t, t)$,
Tweedie's formula recovers an estimate of the clean image:
\begin{equation}\label{eq:tweedie}
  \hat x_0(x_t, t)
  = \frac{x_t - \sqrt{1 - \bar\alpha_t}\, \hat\varepsilon}{\sqrt{\bar\alpha_t}}.
\end{equation}

\medskip
\noindent\textbf{UNet architecture.}
The UNet has 5 resolution levels with channel widths $[64, 128, 256, 512, 512]$,
2 ResNet blocks per level, self-attention at intermediate resolutions,
GroupNorm, and SiLU activations.
Input and output are single-channel $64 \times 64$ images.
The model was pretrained on ImageNet (1.28M grayscale images) and finetuned on
the PNAS natural image dataset (3190 images, 50 epochs).


\subsection{DDIM Sampling}

Rather than running all $T = 1000$ reverse steps (slow), we use DDIM
(Song et al., 2021) which selects a subsequence of $S$ timesteps
$\tau_S > \tau_{S-1} > \cdots > \tau_1$ (we use $S = 50$, uniformly spaced)
and takes deterministic steps:
\begin{equation}\label{eq:ddim}
  x_{\tau_{i-1}}
  = \sqrt{\bar\alpha_{\tau_{i-1}}}\; \hat x_0
  + \sqrt{1 - \bar\alpha_{\tau_{i-1}}}\;
    \underbrace{\frac{x_{\tau_i} - \sqrt{\bar\alpha_{\tau_i}}\, \hat x_0}
    {\sqrt{1 - \bar\alpha_{\tau_i}}}}_{\text{predicted noise direction}},
\end{equation}
where $\hat x_0 = \hat x_0(x_{\tau_i}, \tau_i)$ from Eq.~\eqref{eq:tweedie}.
At the final step ($i = 1$), we set $x_0 = \hat x_0$ directly.

DDIM is deterministic given the initial noise $x_T \sim \N(0, I)$, so different
random seeds produce different images.


%% ============================================================
\section{GP Utility Functions}

We consider two utility functions, both information-theoretic.
Let $R$ denote the (integer) spike count response of the neuron.

\subsection{Standard Utility}

\begin{equation}\label{eq:utility_standard}
  U_{\text{std}}(x^*)
  = H[R \mid x^*, D]
  - H_{\text{noise}}(x^*),
\end{equation}
where $H[R \mid x^*, D]$ is the marginal entropy of the predicted response
and $H_{\text{noise}}$ is the irreducible Poisson noise entropy.
This measures how uncertain we are about the neuron's response at $x^*$
beyond what Poisson noise alone would give.

\subsection{Distribution-Aware (DA) Utility}

\begin{equation}\label{eq:utility_da}
  U_{\text{DA}}(x^*)
  = H[R \mid x^*, D]
  - \E_{\substack{x \sim p(x) \\ \lambda \sim q(\lambda(x) \mid D)}}
    \bigl[ H[R \mid x^*, \lambda(x), D] \bigr].
\end{equation}

The second term averages the \emph{conditional} entropy of $R$ at $x^*$
after hypothetically observing the GP latent $\lambda$ at a random stimulus~$x$
drawn from the image pool.
This accounts for the fact that observing a response at~$x^*$
also reduces uncertainty at nearby points in stimulus space.

\medskip
\noindent\textbf{Entropy computation.}
In both cases, we compute the entropy of $R$ via a Laplace approximation
to the Poisson predictive distribution
$p(r \mid x^*, D) \approx \sum_r P(r) \log P(r)$,
truncated at $r_{\max} = 100$.

\medskip
\noindent\textbf{Differentiability.}
Both utility functions are differentiable with respect to the input image $x^*$
(the gradient flows through the GP kernel computation), which is essential
for the guidance mechanism described next.


%% ============================================================
\section{Guided Reverse Diffusion}

The key idea: at each DDIM denoising step, we compute an estimate of the
clean image $\hat x_0$, evaluate the GP utility $U(\hat x_0)$, and shift
$\hat x_0$ in the direction that increases utility.
This is analogous to classifier guidance (Dhariwal \& Nichol, 2021), but with
the GP utility replacing the classifier.

\subsection{Guidance Gradient}

At step $i$ with timestep $t = \tau_i$, we have the current noisy image $x_t$.
We compute the gradient of utility with respect to $x_t$:
\begin{equation}\label{eq:guidance_grad}
  g_t = \nabla_{x_t} U\!\bigl(\hat x_0(x_t, t)\bigr).
\end{equation}

By the chain rule through Tweedie's formula (Eq.~\ref{eq:tweedie}):
\[
  g_t
  = \frac{\partial \hat x_0}{\partial x_t} \cdot \nabla_{\hat x_0} U(\hat x_0)
  = \frac{1}{\sqrt{\bar\alpha_t}}
    \left(I - \frac{\sqrt{1 - \bar\alpha_t}}{\sqrt{\bar\alpha_t}}
    \frac{\partial \varepsilon_\theta}{\partial x_t}\right)
    \nabla_{\hat x_0} U(\hat x_0).
\]

\medskip
\noindent\textbf{Approximate gradient} (default, faster).
We treat the UNet Jacobian $\partial \varepsilon_\theta / \partial x_t$ as
negligible (run UNet under \texttt{torch.no\_grad}), giving:
\begin{equation}\label{eq:approx_grad}
  g_t^{\text{approx}}
  \approx \frac{1}{\sqrt{\bar\alpha_t}} \nabla_{\hat x_0} U(\hat x_0).
\end{equation}
This halves the compute cost per step.

\subsection{Guided Denoised Estimate}

The utility gradient shifts the Tweedie estimate:
\begin{equation}\label{eq:guided_x0}
  \hat x_0^{\text{guided}}
  = \hat x_0
  + w \cdot \frac{1 - \bar\alpha_t}{\sqrt{\bar\alpha_t}} \cdot g_t,
\end{equation}
where $w$ is the \textbf{guidance scale} (a hyperparameter controlling how
strongly utility influences the generation).

The ramp factor $(1 - \bar\alpha_t) / \sqrt{\bar\alpha_t}$
naturally modulates guidance strength across timesteps:
at high noise ($t \approx T$), the factor is large ($\sim 158$), allowing bold
exploration; at low noise ($t \approx 0$), the factor is small ($\sim 0.3$),
preserving fine image details.

\subsection{Firing Rate Guard}

At high-noise timesteps, $\hat x_0$ can be far from any real image, leading to
extremely high predicted firing rates.
To prevent the gradient from chasing physiologically impossible rates:
\[
  \text{if } \exp(\mu_g) \geq f_{\max}: \quad g_t \leftarrow 0,
\]
where $\mu_g = A \cdot \lambda_{\text{mean}} + \lambda_0$ is the predicted
log-firing rate and $f_{\max} = 100$ Hz.
This guard is typically active only during the first $\sim 10$ DDIM steps.


%% ============================================================
\section{Full Algorithm}

\begin{algorithm}[H]
\caption{Guided Reverse Diffusion for Stimulus Generation}
\label{alg:guided}
\begin{algorithmic}[1]
\Require Trained UNet $\varepsilon_\theta$, noise schedule
  $\{\bar\alpha_t\}_{t=0}^T$, trained GP model,
  DDIM steps $S$, guidance scale $w$, firing rate cap $f_{\max}$,
  scale factor $s$
\State $x_T \sim \N(0, I)$ \Comment{Random initial noise}
\State $\{\tau_S, \tau_{S-1}, \ldots, \tau_1\} \gets$
  uniformly spaced subsequence of $\{0, \ldots, T{-}1\}$
\For{$i = S, S{-}1, \ldots, 1$}
  \State $t \gets \tau_i$
  \Statex
  \Statex \textit{\% --- Step A: Denoise ---}
  \State $\hat\varepsilon \gets \varepsilon_\theta(x_t, t)$
    \Comment{UNet predicts the noise (no gradients through UNet)}
  \State $\hat x_0 \gets
    (x_t - \sqrt{1 - \bar\alpha_t}\, \hat\varepsilon) / \sqrt{\bar\alpha_t}$
    \Comment{Tweedie estimate (Eq.~\ref{eq:tweedie})}
  \Statex
  \Statex \textit{\% --- Step B: Compute utility gradient w.r.t.\ denoised image ---}
  \State $\hat x_0^{\text{raw}} \gets \hat x_0 \cdot s$
    \Comment{Convert to pixel space; must be differentiable}
  \State $U, \mu_g \gets \text{Utility}(\hat x_0^{\text{raw}})$
    \Comment{GP utility at current denoised estimate}
  \State $\nabla_{\hat x_0} U \gets$ gradient of $U$ w.r.t.\ $\hat x_0$
    \Comment{Flows through: GP kernel $\to$ scale factor $\to$ Tweedie}
  \Statex
  \Statex \textit{\% --- Step C: Apply guidance to shift $\hat x_0$ ---}
  \If{$\exp(\mu_g) < f_{\max}$}
    \Comment{Firing rate guard: skip if unrealistic}
    \State $\hat x_0^{\text{guided}} \gets
      \hat x_0 + w \cdot \frac{1 - \bar\alpha_t}{\sqrt{\bar\alpha_t}}
      \cdot \nabla_{\hat x_0} U$
      \Comment{Shift toward higher utility (Eq.~\ref{eq:guided_x0})}
  \Else
    \State $\hat x_0^{\text{guided}} \gets \hat x_0$
    \Comment{No guidance at this step}
  \EndIf
  \Statex
  \Statex \textit{\% --- Step D: DDIM reverse step ---}
  \If{$i > 1$}
    \State $t' \gets \tau_{i-1}$
    \State $d \gets (x_t - \sqrt{\bar\alpha_t}\, \hat x_0^{\text{guided}})
      / \sqrt{1 - \bar\alpha_t}$
      \Comment{Predicted noise direction}
    \State $x_{t'} \gets \sqrt{\bar\alpha_{t'}}\, \hat x_0^{\text{guided}}
      + \sqrt{1 - \bar\alpha_{t'}}\, d$
      \Comment{DDIM step (Eq.~\ref{eq:ddim})}
  \Else
    \State $x_0 \gets \hat x_0^{\text{guided}}$
    \Comment{Final denoised image}
  \EndIf
\EndFor
\State \Return $x_0 \cdot s$ \Comment{Generated image in raw pixel space}
\end{algorithmic}
\end{algorithm}

\noindent\textbf{Note on the gradient in Step~B.}
The utility $U$ is evaluated at $\hat x_0^{\text{raw}}$, which depends on
$\hat x_0$, which depends on $x_t$ via Tweedie's formula
(Eq.~\ref{eq:tweedie}).
In the approximate gradient mode (default), the UNet runs without gradient
tracking, so $\hat\varepsilon$ is treated as a constant.
The gradient $\nabla_{\hat x_0} U$ then flows only through:
\[
  U(\hat x_0^{\text{raw}})
  \;\xleftarrow{\text{chain rule}}\;
  \hat x_0^{\text{raw}} = s \cdot \hat x_0
  \;\xleftarrow{\text{chain rule}}\;
  \hat x_0 = (x_t - c) / \sqrt{\bar\alpha_t},
\]
where $c = \sqrt{1 - \bar\alpha_t}\,\hat\varepsilon$ is a constant (detached).
This yields the approximate gradient from Eq.~\eqref{eq:approx_grad}:
$g_t \approx (1/\sqrt{\bar\alpha_t})\, \nabla_{\hat x_0} U$.
The ramp factor in Step~C absorbs this scaling.

\noindent\textbf{Best-of-$K$ selection.}
Because different random seeds $x_T$ produce different images with varying
utility, we run Algorithm~\ref{alg:guided} with $K$ seeds (typically $K = 100$)
and select the image with the highest final utility.


%% ============================================================
\section{Pixel Space and Physical Constraints}

\subsection{Space Conversion}

The diffusion model operates in a normalized space where pixel values lie
approximately in $[-1, 1]$.
The GP operates in the raw pixel space of the PNAS dataset.
The conversion uses a global scale factor:
\[
  s = \max_{i,j} |x^{\text{PNAS}}_{ij}| = 2.478.
\]
The Tweedie estimate $\hat x_0$ (model space) is converted to raw pixels via
$\hat x_0^{\text{raw}} = s \cdot \hat x_0$.
This multiplication \textbf{must} remain in the computation graph for gradient
flow.

\subsection{Out-of-Bounds Detection}

The physical display (DMD projector) has a finite range.
After generation, we check what fraction of pixels fall outside the dataset's
global $[\min, \max]$ range and flag it as ``OOB'' (out of bounds).
High OOB indicates the guidance pushed pixels beyond physically realizable
values --- the image cannot be shown to the neuron as-is.


%% ============================================================
\section{The Role of the RF Mask}

The arc-cosine kernel used in the GP has a receptive field (RF) structure
controlled by a locality parameter $\beta$.
The kernel applies a Gaussian spatial envelope:
\[
  \alpha_j = \exp\!\left(-\frac{d_j^2}{4\beta^2}\right),
\]
where $d_j$ is the distance of pixel $j$ from the RF center.
Pixels with $\alpha_j < 0.001$ are masked out (set to zero contribution).

\medskip
\noindent\textbf{Effect on guidance.}
Because the kernel ignores pixels outside the mask, the utility gradient
$\nabla_{x^*} U$ is nonzero \emph{only} at RF pixels.
The guidance therefore only modifies pixels within the RF;
the diffusion model controls all remaining pixels freely.

\medskip
\noindent\textbf{Mask size vs.\ $\beta$.}
The mask radius scales linearly with $\beta$:
$r_{\text{mask}} = 2\beta\sqrt{\ln(1/\tau)}$ where $\tau = 0.001$.
The mask \emph{area} (number of included pixels) scales as $\beta^2$.
For a $64 \times 64$ image:

\begin{center}
\begin{tabular}{@{}lcc@{}}
\toprule
$\beta$ & Mask pixels & Coverage \\
\midrule
0.03 & 69 & 1.7\% \\
0.05 & 221 & 5.4\% \\
0.10 (default) & 869 & 21.2\% \\
\bottomrule
\end{tabular}
\end{center}

A smaller mask means fewer pixels receive gradient signal, giving the diffusion
model more control over the final image appearance.
This creates a natural tradeoff:
\begin{itemize}
  \item \textbf{Small $\beta$}: Natural-looking images (low OOB), but limited
    utility gain (few pixels to optimize).
  \item \textbf{Large $\beta$}: High utility achievable, but images become
    distorted (high OOB) as guidance dominates more pixels.
\end{itemize}

\noindent\textbf{Important:} $\beta$ is a learnable kernel parameter.
If not frozen during GP training, the optimizer will change it from the
initial value, and the final mask size will differ from the intended setting.


%% ============================================================
\section{Summary of Key Hyperparameters}

\begin{center}
\begin{tabular}{@{}llp{7cm}@{}}
\toprule
Parameter & Default & Role \\
\midrule
$T$ & 1000 & Total diffusion timesteps \\
$S$ & 50 & DDIM sampling steps \\
$w$ & varies & Guidance scale (how strongly utility influences generation) \\
$K$ & 100 & Number of seeds for best-of-$K$ selection \\
$f_{\max}$ & 100 & Firing rate guard threshold (Hz) \\
$M$ & 50 & Number of GP inducing points \\
$\beta$ & 0.1 & RF mask size (kernel locality parameter) \\
$s$ & 2.478 & Pixel space scale factor (dataset-specific) \\
\bottomrule
\end{tabular}
\end{center}


%% ============================================================
\section*{References}

\begin{itemize}
\item Ho, J., Jain, A., \& Abbeel, P. (2020).
  Denoising Diffusion Probabilistic Models. \textit{NeurIPS}.
\item Song, J., Meng, C., \& Ermon, S. (2021).
  Denoising Diffusion Implicit Models. \textit{ICLR}.
\item Dhariwal, P. \& Nichol, A. (2021).
  Diffusion Models Beat GANs on Image Synthesis. \textit{NeurIPS}.
\end{itemize}

\end{document}
